\section*{Datos de la rapidez de exposición con diferentes materiales}

% Using tabular environment to create a table to simulate the table in the source image.
\begin{table}[h!]
    \begin{tabular}{|c|c|c|c|}
        \hline
        & \textbf{MADERA} & \textbf{ALUMINIO} & \textbf{HIERRO} \\ \hline
        1 PLACA &  &  &  \\ \hline
        2 PLACAS &  &  &  \\ \hline
        3 PLACAS &  &  &  \\ \hline
    \end{tabular}
    \caption{Tabla que muestra los datos de simulacro de la exposición con distintas placas de MADERA, ALUMINIO y HIERRO}
\end{table}

% Creating a multi-chapter figure using minipage to simualte the histograms.
\begin{figure}[h!]
    \centering
    \begin{minipage}{0.3\textwidth}
        \centering
        % TODO
        % \includegraphics[width=\textwidth]{path/to/histogram1.jpg} % replace with actual path
        \caption{mR para MADERA}
    \end{minipage}\hfill
    \begin{minipage}{0.3\textwidth}
        \centering
        % TODO
        % \includegraphics[width=\textwidth]{path/to/histogram2.jpg} % replace with actual path
        \caption{mR para ALUMINIO}
    \end{minipage}\hfill
    \begin{minipage}{0.3\textwidth}
        \centering
        % TODO
        % \includegraphics[width=\textwidth]{path/to/histogram3.jpg} % replace with actual path
        \caption{mR para HIERRO}
    \end{minipage}
\end{figure}

\newpage % Assuming the new page is required as per document structuring.
% No further content was provided, no text or image content continues after this page in provided descriptions. Adjustments or additions can be made above.
